\section{Forward hom-orthogonal sequences}\label{sect2}


We will show that, for any Jacobian algebra $\Lambda=J(Q,W)$ of finite representation type, there is a $1$-$1$ correspondence between maximal green sequences and complete forward hom-orthogonal sequences of Schurian modules, also known as ``bricks''. We prove this by induction on the length of a green sequence. The key step (Lemma~\ref{lem: inductive step A=B for Q finite type}) is known to hold for any cluster-tilted algebra by~\cite{BIRSm}.
But, we go through the details for the benefit of our students.
\subsection{Definition of forward hom-orthogonal sequence}

The results of this subsection hold for $\Lambda$ any finite dimensional algebra over a field.

In the following definition we use the notation $\cF(M):=M^\perp$ for the full subcategory of $\Lambda\text{-}\modr$ of all $Y$ so that $\Hom_\Lambda(M,Y)=0$ and $\cG(M):=\,^\perp\cF(M)$ is the full subcategory of $\Lambda\text{-}\modr$ of all $X$ so that $\Hom_\Lambda(X,Y)=0$ for all $Y\in\cF(M)$. Then $\cF(M)=\cG(M)^\perp$. So, $(\cG(M),\cF(M))$ is a \emph{torsion pair} in $\Lambda\text{-}\modr$ where $\cG(M)$ (\resp  $\cF(M)$) is the \emph{torsion class} (\resp , \emph{torsion-free class}). (See, \eg ,~\cite{BeRe}.) 

The basic properties of the {torsion class} $\cG(M)$ are that $\cG(M)$ is closed under extensions (in particular, direct sums) and also closed under taking quotients. For example, any quotient module $N$ of $M$ lies in $\cG(M)$ and therefore, $\cG(N)\subset\cG(M)$ and $\cF(N)=N^\perp\supset M^\perp=\cF(M)$. We also need the fact that, for any $\Lambda$-module $X$, there is a \emph{canonical short exact sequence}
\[
	0\to rX\to X\to tX\to 0
\]
where $rX\in \cG(M)$ and $tX\in \cF(M)$. The canonical sequence is easy to construct: $rX$ is simply the sum of all submodules of $X$ which are objects of $\cG(M)$. Since $\cG(M)$ is closed under direct sums and quotients, $rX$ is an object in $\cG(M)$.




\begin{prop}\label{prop: HN stratification with 3 strata}
Suppose $\Hom_\Lambda(X,Y)=0$ and $\cC=X^\perp\cap\,^\perp Y$. 
\begin{enumerate}
\item\label{prop2.1_1} If $\cC\neq0$ then $\cC$ contains a Schurian object.
\item\label{prop2.1_2} $\cG(X)=\,^\perp \cC\cap\,^\perp Y$.
\end{enumerate}
\end{prop}

\begin{proof} \ \\*[-0.8em]

 \eqref{prop2.1_1} Let $M$ be a nonzero object of $\cC$ of minimal length. Then we claim that $M$ is Schurian. If not, there would be an endomorphism $f:M\to M$ whose image $f(M)$ is nonzero and of smaller length than $M$. But $f(M)\in X^\perp$, being a submodule of $M$ and $f(M)\in \,^\perp Y$ since it is a quotient object of $M$. Therefore, $f(M)$ lies in $\cC$ which contradicts the minimality of $M$.

\eqref{prop2.1_2} $\cG(X)\subset \,^\perp Y$ since $Y\in \cF(X)$. $\cG(X)\subset \,^\perp \cC$ since $\cC\subset  X^\perp=\cF(X)$. Therefore, $\cG(X)\subset \,^\perp\cC\cap \,^\perp Y$. Conversely, let $M\in \,^\perp \cC\cap\,^\perp Y$. Then, in the canonical short exact sequence $0\to rM\to M\to tM\to 0$ for the torsion pair $(\cG(X),\cF(X))$, $tM\in  \,^\perp \cC\cap\,^\perp Y$ since $tM$ is a quotient of $M\in \,^\perp \cC\cap\,^\perp Y$. Also, $tM\in \cF(X)= X^\perp$ by definition of the canonical sequence. So, $tM\in X^\perp\cap \,^\perp\cC\cap\,^\perp Y=\cC\cap \,^\perp \cC=0$. So, $M=rM\in\cG(X)$.
\end{proof}

\begin{defi}\label{def of FHO}
A \emph{forward hom-orthogonal (FHO) sequence} in $\Lambda\text{-}\modr$ is a finite sequence of Schurian modules $M_1, \cdots, M_m$ so that
\begin{enumerate}
\item\label{defi2.2_1} $\Hom_\Lambda(M_i,M_j)=0$ for all $1\le i<j\le m$.
\item\label{defi2.2_2} The sequence is maximal in $\cG(M)$ where $M=M_1\oplus \cdots\oplus M_m$. By \emph{maximal} we mean that no other Schurian object in $\cG(M)$ can be inserted into the sequence $M_1,\ldots,M_m$ preserving property~\eqref{defi2.2_1}.
\end{enumerate}
We say that a forward hom-orthogonal sequence is \emph{complete} if $\cG(M)=\Lambda\text{-}\modr$ or, equivalently, $\cF(M)=0$ where we generally let $M$ denote $M_1\oplus\cdots\oplus M_m$ when it is clear from the context which forward hom-orthogonal sequence is being considered. If the sequence $(M_i)$ satisfies only~\eqref{defi2.2_1} we call it \emph{weakly forward hom-orthogonal}. Note that, by Proposition~\ref{prop: HN stratification with 3 strata}, property~\eqref{prop2.1_1} is equivalent to the condition that $\cG(M)\cap(M_1\oplus \cdots \oplus M_k)^\perp\cap\,^\perp(M_{k+1}\oplus\cdots\oplus M_m)=0$ for all $k$ since, if this intersection were nonzero, by~\ref{prop: HN stratification with 3 strata}~\eqref{prop2.1_1} it would contain a Schurian object which could be inserted between $M_k$ and $M_{k+1}$ in the FHO sequence.
\end{defi}

As a consequence of Proposition~\ref{prop: HN stratification with 3 strata} we have the following.


\begin{coro}\label{cor: max weakly forward hom-orth is complete}
Let $M_1,\ldots,M_m$ be a weakly FHO sequence of Schurian $\Lambda$-modules which is \emph{maximal} in the sense that it is not a proper subsequence of any other weakly FHO sequence. Then $M_1,\ldots,M_m$ is complete.
\end{coro}

\begin{proof}
It follows from Definition~\ref{def of FHO} that $M_1,\ldots,M_m$ is FHO in $\cG(M)$. It remains to show that $\cG(M)=\Lambda\text{-}\modr $. Suppose not. Then, $\cF(M)\neq0$ and, by Proposition~\ref{prop: HN stratification with 3 strata} with $Y=0$ and $X=M$, there would exist a Schurian module in $\cF(M)=M^\perp$. This Schurian module could be added after $M_m$ contradicting its maximality.
\end{proof}

\begin{prop}\label{prop:subsequence of hom orth is hom orth}
If $M_1,\ldots,M_m$ is a FHO sequence in $\Lambda\text{-}\modr $ then so is $M_1,\ldots,M_k$ for any $k<m$.
\end{prop}

\begin{proof}
Suppose not. Then there is a Schurian object $X$ in $\cG(M_1\oplus \cdots\oplus M_k)$ which can be inserted into the sequence $M_1,\ldots,M_k$. But $X\in \,^\perp M_j$ for $k<j\le m$ since $X\in \cG(M_1\oplus\cdots\oplus M_k)\subset \cG(M_1\oplus\cdots\oplus M_m)$. So, $X$ can also be inserted into the original sequence $M_1,\ldots,M_m$ giving a contradiction.
\end{proof}

\begin{lemm}\label{lem:when Sk is not one of the Mi}
Let $M_1,\ldots,M_m$ be a FHO sequence in $\Lambda\text{-}\modr $ so that no $M_i$ is isomorphic to the simple module $S_j$. Then $\Hom_\Lambda(M_i,S_j)=0$ for all $i$.
\end{lemm}

\begin{proof}
Suppose not and let $i$ be minimal so that $\Hom_\Lambda(M_i,S_j)\neq0$. Since $S_j$ is simple, any nonzero morphism $M_i\to S_j$ must be an epimorphism. So, $S_j\in\cG(M)$, where $M=\bigoplus_iM_i$. Also, $\Hom_\Lambda(S_j,M_k)=0$ for all $i\le k\le m$. So, $S_j$ can be inserted into the FHO sequence to the left of $M_i$ contradicting its maximality.
\end{proof}

This Lemma implies that, if a simple module is missing from a FHO sequence, it can be added after the last module.

\begin{prop}
Each complete FHO sequence in $\Lambda\text{-}\modr $ contains every simple $\Lambda$-module.\qed
\end{prop}

\begin{coro}\label{cor: every comple fhos ends in a simple}
If $M_1,\ldots,M_m$ is a complete FHO sequence then $M_m$ is simple.
\end{coro}

\begin{proof}
If $M_m$ were not simple, all simple $\Lambda$-modules would come before $M_m$ in the sequence. Then, $\Hom_\Lambda(S_k,M_m)=0$ for all $k$ which is impossible.
\end{proof}


\begin{lemm}\label{lem:FHO of length one is simple}
Let $M$ be a Schurian module. Then $M=M_1$ is a FHO sequence of length $1$ if and only if $M$ is simple.
\end{lemm}

\begin{proof}
Let $M=S_i$ be the $i$th simple module. Then $\cF(S_i)$ contains the injective modules $I_j$ for all $j\neq i$. So, $\cG(S_i)$ consists only of iterated self extensions of $S_i$. So, $S_i$ is the only Schurian object in $\cG(S_i)$. So, it is a FHO sequence of length 1.

Conversely, let $M=M_1$ be a FHO sequence of length 1. If $M$ is not simple then, by Lemma~\ref{lem:when Sk is not one of the Mi}, $\Hom_\Lambda(M,S_j)=0$ for all $S_j$ which is impossible.%Let $S$ be a simple quotient of $M_1$. Since $\cG(M)$ is closed under quotient objects, $S\in\cG(M)$. If $M_1\neq S$, then $\Hom_\Lambda(S,M_1)=0$ since $M_1$ is Schurian. So, $S,M_1$ weakly FHO in $\cG(M)$ contradicting the maximality of $M$. So, $M_1=S$ is simple.
\end{proof}

Proposition~\ref{prop:subsequence of hom orth is hom orth}, Lemma~\ref{lem:FHO of length one is simple} and Proposition~\ref{cor: every comple fhos ends in a simple} imply that every FHO sequence starts with a simple module. This leads to the following very useful result.

\begin{prop}\label{prop: complete iff M1 is simple and M2, etc is maximal}
Let $M_1,\ldots,M_m$ be a sequence of Schurian $\Lambda$-modules. Then the following are equivalent.
\begin{enumerate}[(1)]
\item\label{prop2.9_1} $M_1,\ldots,M_m$ is a complete FHO sequence in $\Lambda\text{-}\modr $.

\item\label{prop2.9_2} $M_1$ is simple and $M_2,\ldots,M_m$ is a maximal weakly FHO sequence in $M_1^\perp$.

\item\label{prop2.9_3} $M_m$ is simple and $M_1,\ldots,M_{m-1}$ is a maximal weakly FHO sequence in $\,^\perp M_m$.
\end{enumerate}
\end{prop}

\begin{proof} Since every complete FHO sequence starts and ends with a simple module, \ref{prop2.9_1}~implies~\ref{prop2.9_2} and~\ref{prop2.9_1}.

For the converse, we observe in both cases~\ref{prop2.9_2} and~\ref{prop2.9_3}, every simple module must occur in the sequence. Suppose not. Let $S_k$ be a simple module not isormorphic to any $M_i$. In Case~\ref{prop2.9_3}, $M_1,\ldots,M_{m-1}$ is FHO by definition since $\,^\perp M_m$ is a torsion class. Therefore, by Lemma~\ref{lem:when Sk is not one of the Mi}, $\Hom_\Lambda(M_i,S_k)=0$ for all $i<m$. Also, $\Hom_\Lambda(S_k,M_m)=0$ since $S_k,M_m$ are nonisomorphic simple modules. So, $S_k$ can be added to the sequence $M_1,\ldots,M_{m-1}$ contradicting its maximality as a sequence in $\,^\perp M_m$. In case (2), let $j\ge1$ be maximal so that $\Hom_\Lambda(M_i,S_k)=0$ for $1\le i\le j$. Then, either $j<m$ in which case $\Hom_\Lambda(M_{j+1},S_k)\neq0$ or $j=m$. In either case, $S_k$ can be inserted into the sequence $M_1,\ldots,M_m$ after $M_j$, where $j\ge1$, contradicting the maximality of $M_2,\ldots,M_m$ in $M_1^\perp$. 

This implies that $M_1,\ldots,M_m$ is maximal weakly FHO and thus complete (by Corollary~\ref{cor: max weakly forward hom-orth is complete}). Indeed, in Case~\ref{prop2.9_2}, no module can be inserted before $M_1$ since every module maps to at least one simple module. No Schurian module can be inserted after $M_1$ since that would increase the length of the sequence $M_2,\ldots,M_m$ in $M_1^\perp$ contradicting its maximality. Case~\ref{prop2.9_3} is similar. Therefore, the three conditions are equivalent.
\end{proof}


\subsection{Rotation of forward hom-orthogonal sequences} 


We assume for the rest of the paper that $Q$ is a quiver without loops or oriented 2-cycles which is of finite type in the sense of Fomin and Zelevinsky~\cite{FZ4}. Let $W$ be a nondegenerate potential for $Q$ in the sense of~\cite{DWZ1}. Let $\Lambda=J(Q,W)$ be the Jacobian algebra of the quiver with potential $(Q,W)$. Mutation of quivers with potential will be reviewed below (before the proof of Lemma~\ref{lem: inductive step A=B for Q finite type}).

We will show, in Theorem~\ref{thm: inductive thm implying main thm}, that sequences of $m$ green mutations (i.e., mutations on positive $c$-vectors) of the quiver $Q$ are in bijection with FHO sequences of length $m$ in $J(Q,W)\text{-}\modr $. This will immediately imply the bijection (Corollary~\ref{cor: main thm of section 2}):
\[
	\left\{\begin{array}{cc}
	\text{maximal green sequences}\\
	\text{ for $Q$}
	\end{array}
	\right\}
\cong 
	\left\{\begin{array}{cc}
	\text{complete FHO sequences}\\
	\text{ for $J(Q,W)$}
	\end{array}
	\right\}.
\]



We will prove the $m$ mutation statement by induction on $m$ using Lemma~\ref{lem: inductive step A=B for Q finite type} below which, by Corollary~\ref{cor: rotation lemma}, is an analogue of the Rotation Lemma~\cite{BHIT} for complete FHO sequences. Before this we review the mutation process for quivers with potential: $(Q',W')=\mu_k(Q,W)$ assuming that $Q$ is a quiver of finite type with nondegenerate potential $W$. We begin with the observation that, since $Q$ has finite type, it cannot contain the subquiver 
\[
	\tilde A_2:\xymatrix{ i\ar@/^1pc/[rr]\ar[r] &k\ar[r] & j}
\]
since, if it did, mutation at $k$ would produce a double arrow from $i$ to $j$.


Let $I$ be the set of all vertices $i$ for which there is an arrow $\alpha_i:i\to k$ in $Q$. Let $J$ be the set of all vertices $j$ for which there is an arrow $\beta_j:k\to j$ in $Q$. Since $Q$ has finite type, $Q$ has no arrows between elements of $I$ and no arrows between elements of $J$. Otherwise $Q$ would contain a subquiver of type $\tilde A_2$. Additionally, for each $(i, j) \in  I \times J$ there are no arrows $i \rightarrow j$ and there is at most one arrow $j \rightarrow i$.



Let $P$ be the set of all pairs $(i,j)\in I\times J$ for which $Q$ has an arrow $\gamma_{ij}:j\to i$. Let $P'$ be the complement of $P$ in $I\times J$. The mutated quiver $Q'=\mu_kQ$ is obtained from $Q$ by reversing the orientations of the arrows $\alpha_i,\beta_j$ to produce new arrows $\alpha_i^\ast:k\to i$, $\beta_j^\ast:j\to k$, adding an arrow $\gamma_{i'j'}^\ast:i'\to j'$ for each $(i',j')\in P'$ and deleting the arrows $\gamma_{ij}:j\to i$ for all $(i,j)\in P$. This is indicated as follows.% where $B_{ij}, 
\[
%\xymatrixrowsep{10pt}\xymatrixcolsep{10pt}
\xymatrix{%begin xy matrix
Q: & j\ar@/^1pc/[rr]^{\gamma_{ij}} &k\ar[l]^{\beta_j} & i\ar[l]^{\alpha_i}%\ar@{--}[ld]
&& Q'=\mu_k Q: & j'\ar[r]_{\beta_{j'}^\ast}%\ar@{--}[rd] 
&k \ar[r]_{\alpha_{i'}^\ast} & i'\ar@/_1pc/[ll]_{\gamma_{i'j'}^\ast}
	}%end xy matrix
\]

The potential $W$ is a linear combination of oriented cycles in $Q$. We begin with the general form of such a potential then simplify it by a change of coordinates. Collecting together all terms involving $\gamma_{ij}$ or $\beta_j\alpha_i$ for $(i,j)\in P$, we have:
\[
	W=\sum_{(i,j)\in P}a_{ij}\gamma_{ij}\beta_j\alpha_i+\sum_{(i,j)\in P}A_{ij}\beta_j\alpha_i+\sum_{(i,j)\in P}\gamma_{ij}B_{ij}+ \text{ other terms}
\]
where $a_{ij}$ are nonzero scalars, $A_{ij},B_{ij}$ are linear combinations of path $j\to i$ and $i\to j$ of length at least 2. We claim that, by changing the choice of $\gamma_{ij}$ we may eliminate the terms $A_{ij}$. The new $\gamma_{ij}$ will be equal to $a_{ij}\gamma_{ij}$ plus a converging, possibly infinite sum of paths $j\to i$.

The first step in this simplification process is given by changing the choice of $\gamma_{ij}$ to $\gamma_{ij}':=a_{ij}\gamma_{ij}+A_{ij}$. Equivalently, we make the substitution 
\[
\gamma_{ij}=a_{ij}^{-1}\gamma_{ij}'-a_{ij}^{-1}A_{ij}.
\]
Then we obtain:
\[
	W=\sum_{(i,j)\in P}\gamma_{ij}'\beta_j\alpha_i+\sum_{(i,j)\in P}A_{ij}'\beta_j\alpha_i+\sum_{(i,j)\in P}\gamma_{ij}'B_{ij}'+ \text{ other terms}.
\]
The old terms $A_{ij}\beta_j\alpha_i$ are cancelled. However, new term arise from $\gamma_{ij}B_{ij}=a_{ij}^{-1}\gamma_{ij}'B_{ij}-a_{ij}^{-1}A_{ij}B_{ij}$. The term $a_{ij}^{-1}A_{ij}B_{ij}$ might give a term $A_{ij}'\beta_j\alpha_i$ since $\beta_j\alpha_i$ might be a subword of a term in $A_{ij}$. To fix this we let $\gamma_{ij}'':=\gamma_{ij}'+A_{ij}'$ and substitute again:
\[
	\gamma_{ij}'=\gamma_{ij}''-A_{ij}'.
\]
This process converges because the minimum length of the paths in the sum $A_{ij}$ is going to infinity. To see this, define the ``adjusted length'' of each path in $A_{ij}$ to be its length minus the number of two letter subwords of the form $\beta_j\alpha_i$. When a cycle in the sum $A_{ij}B_{ij}$ becomes a cycle in $A_{ij}'\beta_j\alpha_i$, the length of the path in $A_{ij}'$ is at least equal to the length of the corresponding path in $A_{ij}$ since $B_{ij}$ is a linear combination of paths of length $\ge2$. However, the adjusted length strictly increases since the subword $\beta_j\alpha_i$ has been removed from $A_{ij}$. Therefore, the process converges. So, we may assume $a_{ij}=1$ and $A_{ij}=0$ in the expression for $W$ and we have:
\[
	W=\sum_{(i,j)\in P}\gamma_{ij}\beta_j\alpha_i+F
\]
where $F$ is a linear combination of oriented cycles in $Q$ none of which contain the subword $\beta_j\alpha_i$ for any $(i,j)\in P$. We may also assume that each of these oriented cycles, written as a path, does not start at vertex $k$. Then, if $\alpha_i$ occurs as a letter in one of these paths, the next letter must be $\beta_{j}$ where $(i,j)\in P'$. This implies that, for each $i\in I$, we have
\[
	\partial_{\alpha_i}F=\sum_{j:(i,j)\in P'} F_{ij}\beta_j
\]
where $F_{ij}$ is a linear combination of paths from $j$ to $i$. For each occurrence of the word $\beta_j\alpha_i$ in the terms in $F$ we get one term in $F_{ij}$. With this description we also see~that
\[
	\partial_{\beta_j}F=\sum_{i:(i,j)\in P'} \alpha_iF_{ij}.
\]
For each $(i,j)\in P$ let
\[
	G_{ij}:=\partial_{\gamma_{ij}}F.
\]


The Jacobian algebra $J(Q,W)$ is the path algebra of $Q$ modulo the following relations.
\[
	(\forall (i,j)\in P)\quad \partial_{\gamma_{ij}}W=0:\qquad\qquad \beta_j\alpha_i=-G_{ij}\qquad\quad 
\]
\begin{equation}\label{partial alpha-i W=0}
	(\forall i\in I)\quad \partial_{\alpha_i}W=0:\qquad \sum_{j:(i,j)\in P} \gamma_{ij}\beta_j+\sum_{j':(i,j')\in P'} F_{ij'}\beta_{j'} =0
\end{equation}
\[
	(\forall j\in J)\quad \partial_{\beta_j}W=0:\qquad \sum_{i: (i,j)\in P} \alpha_i\gamma_{ij}+\sum_{i': (i',j)\in P'} \alpha_{i'}F_{i'j} =0
\]
And, for all other arrows $\delta$, the equation $\partial_\delta W=0$ is just $W_\delta:=\partial_\delta F=0$.

We need the following notation. Let $X$ be a linear combination of paths or oriented cycles in $Q$ which do not start or end at vertex $k$ and which do not contain the subword $\beta_j\alpha_i$ for any $(i,j)\in P$, for example, $X=F,F_{ij},G_{ij}$. We denote by $\widetilde X^\ast$ the corresponding linear combination of paths or cycles in $Q'$ given by replacing each occurrence of the letter $\gamma_{ij}$ in $X$, for $(i,j)\in P$, with $\alpha_i^\ast\beta_j^\ast$ and each occurrence of the pair of letters $\beta_{j'}\alpha_{i'}$, for $(i',j')\in P'$, with $\gamma_{i'j'}^\ast$. 

We need formulas for commuting cyclic derivatives with the operation $\widetilde{(\ )}^\ast$. The best explanation might be by example: Given an oriented cycle 
\[X=(\gamma_{ij})ab(\beta_{j'}\alpha_{i'})de(\gamma_{ij'})gh
\]
 we have $\widetilde X^\ast=
(\alpha_i^\ast\beta_j^\ast)  ab( \gamma_{i'j'}^\ast )de(\alpha_i^\ast\beta_{j'}^\ast )gh$. Then 
\[
\partial_{\alpha_i^\ast} \widetilde X^\ast=\beta_j^\ast  ab( \gamma_{i'j'}^\ast )de(\alpha_i^\ast\beta_{j'}^\ast )gh+
\beta_{j'}^\ast gh (\alpha_i^\ast\beta_j^\ast)  ab( \gamma_{i'j'}^\ast )de
= \beta_j^\ast \widetilde{\partial_{\gamma_{ij}}X}^\ast+\beta_{j'}^\ast \widetilde{\partial_{\gamma_{ij'}}X}^\ast
\]
following the general formula:
\[
	\partial_{\alpha_i^\ast} \widetilde X^\ast=\sum_{j':(i,j')\in P'}\beta_{j'}^\ast \widetilde{\partial_{\gamma_{ij'}}X}^\ast .
\]
Using this formula, we obtain the following calculations.
\begin{equation}\label{eq: derivative of tilde F}
	\forall i\in I:\qquad \partial_{\alpha_i^\ast}\widetilde F^\ast=\sum_{j:(i,j)\in P}\beta_{j}^\ast \widetilde{\partial_{\gamma_{ij}}F}^\ast
	=\sum_{j:(i,j)\in P}\beta_{j}^\ast \widetilde G_{ij}^\ast
\end{equation}
\[
	\forall j\in J:\qquad \partial_{\beta_j^\ast}\widetilde F^\ast=\sum_{i:(i,j)\in P} \widetilde G_{ij}^\ast\alpha_{i}^\ast
\]
\[
	\forall (i',j')\in P': \qquad \partial_{\gamma_{ij}^\ast}\widetilde F^\ast=\widetilde F_{ij}^\ast
\]
All other arrows $\delta$ in $Q'$ are also arrows in $Q$ and the operations $\partial_\delta$ and $\widetilde{(\ )}^\ast$ commute:
\[
	\partial_\delta \widetilde F^\ast=\widetilde{\partial_\delta F}^\ast.
\]


\begin{prop}\label{prop: formula for mutation of W}  
The mutation $\mu_k(Q,W)$ of this quiver with potential in direction $k$ is equivalent to $(Q',W')$ where $Q'$ is defined above and
\[
	W'=-\sum_{(i',j')\in P'} \alpha_{i'}^\ast\beta_{j'}^\ast\gamma_{i'j'}^\ast+\widetilde F^\ast
\]
where the notation $\widetilde F^\ast$ is defined above. The Jacobian algebra $J(Q',W')$ is the path algebra of $Q'$ modulo the relations given as follows.
\[
	(\forall i\in I)\quad \partial_{\alpha_i^\ast}W'=0:\qquad -\sum_{j':(i,j')\in P'} \beta_{j'}^\ast\gamma_{ij'}^\ast +
	\sum_{j:(i,j)\in P}\beta_{j}^\ast \widetilde G_{ij}^\ast=0
\]
\[
	(\forall j\in J)\quad \partial_{\beta_j^\ast}W'=0:\qquad -\sum_{i': (i',j)\in P'} \gamma_{i'j}^\ast\alpha_{i'}^\ast+\sum_{i:(i,j)\in P} \widetilde G_{ij}^\ast\alpha_{i}^\ast =0
\]
\[
	(\forall (i',j')\in P')\quad \partial_{\gamma_{i'j'}^\ast}W'=0:\qquad\qquad \alpha_{i'}^\ast\beta_{j'}^\ast=\widetilde F_{i'j'}^\ast\qquad\quad 
\]
And, for all other arrows $\delta$, the equation $\partial_\delta W'=0$ is just $\widetilde W_\delta^\ast=0$ where $W_\delta=\partial_\delta  F$.
\end{prop}

\begin{proof}
We follow the original definition of mutation of a quiver with potential as given in~\cite{DWZ1}. The first step is to replace the pair of letters $\beta_j\alpha_i$ with a new letter $\gamma_{ij}^\ast$ for all $(i,j)\in I\times J=P\coprod P'$. We denote the result of such a procedure with a tilde $\widetilde\,$. We~get:
\[
	\widetilde W= \sum_{(i,j)\in P} \gamma_{ij}\gamma_{ij}^\ast + \widetilde F
\]
where $\widetilde F$ is a linear combination of oriented cycles of length at least 3. The next step is usually to add new terms $\alpha_i^\ast\beta_j^\ast\gamma_{ij}^\ast$ for all $(i,j)\in I\times J$. However, it is more convenient to subtract these terms to get:
\[
	\widetilde W^+=-\sum_{(i,j)\in I\times J} \alpha_i^\ast\beta_j^\ast\gamma_{ij}^\ast+ \sum_{(i,j)\in P} \gamma_{ij}\gamma_{ij}^\ast + \widetilde F
\]
This deviation from standard procedure is justified since the letters $\alpha_i^\ast$ occur only in this new sum. By replacing each $\alpha_i^\ast$ with its negative, we change the signs of these new terms.

The next step is to change the basis again by replacing $\gamma_{ij}$ with $\gamma_{ij}'+\alpha_i^\ast\beta_j^\ast$ for each $(i,j)\in P$. This eliminates those terms in the first sum corresponding to $(i,j)\in P$ to~give:
\[
	\widetilde W'=-\sum_{(i,j)\in P'} \alpha_i^\ast\beta_j^\ast\gamma_{ij}^\ast+ \sum_{(i,j)\in P} \gamma_{ij}'\gamma_{ij}^\ast + \widetilde F'
\]
The term $\widetilde F'$ contains the term that we want $\widetilde F^\ast$ plus other unwanted terms involving $\gamma_{ij}'$ for $(i,j)\in P$. So, $\widetilde F'=\widetilde F^\ast+\sum \gamma_{ij}'Z_{ij}$ where $Z_{ij}$ is a linear combination of paths of length at least 2 from $i$ to $j$ for each $(i,j)\in P$. Then $\widetilde W'$ can be rewritten as:
\[
	\widetilde W'=-\sum_{(i,j)\in P'} \alpha_i^\ast\beta_j^\ast\gamma_{ij}^\ast+ \sum_{(i,j)\in P} \gamma_{ij}'\gamma_{ij}^\ast+\sum_{(i,j)\in P} \gamma_{ij}'Z_{ij} + \widetilde F^\ast.
\]
Since $\gamma_{ij}^\ast$ only occurs in the second sum, we can change coordinates, replacing $\gamma_{ij}^\ast$ with $\gamma_{ij}'^{\ast}-Z_{ij}$ to eliminate the third, unwanted sum. Then
\[
	\widetilde W'=-\sum_{(i,j)\in P'} \alpha_i^\ast\beta_j^\ast\gamma_{ij}^\ast+ \sum_{(i,j)\in P} \gamma_{ij}'\gamma_{ij}'^\ast+ \widetilde F^\ast.
\]
Now the letters $\gamma_{ij}'$ and $\gamma_{ij}'^\ast$ occur only in the second sum which is a sum of 2-cycles. The final step is to ``reduce'' this expression by deleting these isolated 2-cycles to give
\[
	 W'=-\sum_{(i,j)\in P'} \alpha_i^\ast\beta_j^\ast\gamma_{ij}^\ast+  \widetilde F^\ast
\]
as claimed. Computation of the cyclic derivatives of $W'$ is given by Formulas~\ref{eq: derivative of tilde F}.
\end{proof}



%Using the above notation for $(Q',W')=\mu_k(Q,W)$ we give the proof of Lemma~\ref{lem: inductive step A=B for Q finite type}.


The following lemma is proved in much greater generality in~\cite[Sec~7]{BIRSm} where the statement is demonstrated for any cluster-tilted algebra of the form $\Lambda=J(Q,W)$. The statement uses an involution $\varphi_k$ on $\ZZ^n$ defined as follows.

For any $x \Mk \in \Mk \ZZ^n$ and any $k\Mk \in \Mk Q_0$, let $\varphi_k(x)\Mk =\Mk y\in \ZZ^n$ be given by $y_i\Mk =\Mk x_i$ for $i\Mk \neq \Mk k$~and
\begin{equation}\label{eq: y=Xkx}
	y_k=-x_k+\sum_{k\to j} x_j
\end{equation}
where the sum is over all arrows $k\to j$ in $Q$.
Since $Q$ has no loops, $\varphi_k$ is an involution: $x=\varphi_k(y)$.


\begin{lemm}\label{lem: inductive step A=B for Q finite type}
Let $S_k$ be a simple $\Lambda$-module for $\Lambda=J(Q,W)$ and let $\Lambda'$ be the Jacobian algebra of $(Q',W')=\mu_k(Q,W)$ where $W$ is a nondegenerate potential for $Q$. Suppose that $Q$ has finite type. Then there is an equivalence of full subcategories $\psi_k:S_k^\perp\cong \,^{\perp}S_k'$ where
\[
\begin{aligned}
	S_k^\perp&=\{X\in \Lambda\text{-}\modr \mid %\,|\, 
	\Hom_\Lambda(S_k,X)=0\}
\\
	\,^{\perp}S_k'&=\{Y\in \Lambda'\text{-}\modr\mid \Hom_{\Lambda'}(Y,S_k')=0\}
\end{aligned}
\]
with $S_k,S_k'$ being the simple $\Lambda,\Lambda'$-modules at vertex $k$. Furthermore, the dimension vectors of $X$ and $Y=\psi_k(X)$ are related by
\[
	\undim \psi_k(X)=\varphi_k(\undim X).
\]
where $\varphi_k$ is the automorphism of $\ZZ^n$ given by~\eqref{eq: y=Xkx} above.
\end{lemm}


\begin{proof}%[Proof of Lemma~\ref{lem: inductive step A=B for Q finite type}]



Any representation $X$ of $J(Q,W)$ is given by vector spaces $X_v$ for all vertices $v$ of $Q$ and linear maps $X_v\to X_w$ for arrows $v\to w$ satisfying the relations $\partial_\alpha W=0$. In particular, the arrows $k\to j$, for $j\in J$, induce linear maps $\beta_j:X_k\to X_j$. If $X\in S_k^\perp$, the sum of these linear maps will give a monomorphism $(\beta_j):X_k\to \bigoplus X_j$. Let $Y_k$ with maps $\beta_j^\ast:X_j\to Y_k$ be the cokernel. This gives a functorial short exact sequence:
\begin{equation}\label{functorial exact sequence}
	0\to X_k\xrightarrow{(\beta_j)}\bigoplus_{j\in J} X_j\xrightarrow{(\beta_j^\ast)}Y_k\to 0.
\end{equation}

Let $Y=\psi_k(X)$ be the representation of $\Lambda'=J(Q',W')$ given as follows. For each vertex $s\neq k$, let $Y_s=X_s$. $Y_k$ is given in~\eqref{functorial exact sequence}. Then the exact sequence~\eqref{functorial exact sequence} immediately implies
\[
	\undim Y=\varphi_k(\undim X)
\]
where $\varphi_k$ is the automorphism of $\ZZ^n$ given by~\eqref{eq: y=Xkx} above. For each $j\in J$, the map $\beta_j^\ast:Y_j\to Y_k$ is given in~\eqref{functorial exact sequence} since $Y_j=X_j$. Since these give an epimorphism $\bigoplus Y_j\onto Y_k$, we will have $\Hom_{\Lambda'}(Y,S_k)=0$. So $Y\in\,^{\perp}S_k'$ (once we show that $Y$ is a representation of $J(Q',W')$). For each $(i',j')\in P'$, let $\gamma_{i'j'}^\ast= \beta_{j'}\alpha_{i'}: Y_{i'}\to Y_{j'}$. If $\delta:s\to t$ is an arrow in $Q'$ which is not equal to $\alpha_i^\ast,\beta_j^\ast$ or $\gamma_{ij}^\ast$ then $\delta$ is also an arrow in $Q$ and we define $\delta:Y_s\to Y_t$ to be the equal to the morphism $\delta:X_s\to X_t$. It remains to define $\alpha_i^\ast:Y_k\to Y_i$ for all $i\in I$.


Given $i\in I$, consider the morphism ${(\gamma_{ij},F_{ij'})}:\bigoplus X_j\to X_i$ given by $\gamma_{ij}$ on $Y_j$ for all $j$ so that $(i,j)\in P$ and by $F_{ij'}$ on $Y_{j'}$ for all $j'$ so that $(i,j')\in P'$. Then, Equation~\eqref{partial alpha-i W=0} implies that the composition ${(\gamma_{ij},F_{ij'})}\circ (\beta_j)=0$. Therefore, there is a unique induced map $\alpha_i^\ast:Y_k\to Y_i=X_i$ satisfying the following. (See Equation~\eqref{definition of alpha i ast}.)
\begin{enumerate}
\item\label{proof_lemma2.11_1} $\alpha_i^\ast \beta_j^\ast=\gamma_{ij}$ when $(i,j)\in P$,
\item\label{proof_lemma2.11_2} $\alpha_i^\ast \beta_{j'}^\ast=F_{ij'}$ when $(i,j')\in P'$.
\end{enumerate}
\begin{equation}\label{definition of alpha i ast}
%\xymatrixrowsep{10pt}\xymatrixcolsep{10pt}
\xymatrix{%begin xy matrix
0 \ar[r] &
	X_k\ar@/_1pc/[rd]_(.35)0\ar[r]^{(\beta_j)} &
	\bigoplus X_j\ar[d]_(.4){(\gamma_{ij},F_{ij'})}\ar[r]^{(\beta_j^\ast)}&
Y_k \ar[r]\ar@{-->}[dl]^(.4){\alpha_i^\ast}& 0\\
 &  & X_i=Y_i &  
	}%end xy matrix
\end{equation}

We need to verify that $Y$, with these maps, satisfies the relations for $J(Q',W')$. Note that since $\alpha_i^\ast \beta_j^\ast=\gamma_{ij}$ for any $(i,j)\in P$ and $\gamma_{i'j'}^\ast= \beta_{j'}\alpha_{i'}$ for $(i',j')\in P'$, the map $\lambda:X_s\to X_t$ for any path $\lambda$ in $Q$ with $s,t\neq k$ which does not contain a subpath $\alpha_i \beta_j$ for any $(i,j)\in P$ is equal to the map $\widetilde\lambda^\ast:Y_s\to Y_t$ for $Y$:
\[
	\widetilde\lambda^\ast=\lambda:X_s=Y_s\to X_t=Y_t.
\]
In particular, this implies the required condition that $\widetilde W_\delta^\ast=0$ for $Y$ for any arrow $\delta$ not equal to $\alpha_i^\ast,\beta_j^\ast$ or $\gamma_{i'j'}^\ast$ since $\widetilde W_\delta^\ast=W_\delta=0$ on $X$ and thus on $Y$.

Condition (2) above is equivalent the required condition $\alpha_{i'}^\ast \beta_{j'}^\ast=\widetilde F^\ast_{i'j'}$ for $(i',j')\in P'$ ($\partial_{\gamma_{i'j'}^\ast}W'=0$) since $\widetilde F^\ast_{i'j'}= F_{i'j'}$ on $Y$.

To verify the relation $\partial_{\alpha_i^\ast}W'=0$ we substitute $\gamma_{ij'}^\ast=\beta_{j'}\alpha_i$ for $(i,j')\in P'$ and $\widetilde G_{ij}^\ast=G_{ij}=-\beta_j\alpha_i$ for $(i,j)\in P$ (from~\eqref{partial alpha-i W=0}). Then, the required condition $\partial_{\alpha_i^\ast}W'=0$ becomes:
\[
	\sum_{(i,j')\in P'}\beta_{j'}^\ast\beta_{j'}\alpha_i+\sum_{(i,j)\in P}\beta_j^\ast\beta_j\alpha_i=0
\]
But $P\coprod P'=I\times J$. So, this is the same as $\sum_{ij} \beta_j^\ast\beta_j\alpha_i=0$ which follows from the fact that $\sum_j\beta_j^\ast\beta_j=0$.

Similarly, the relation $\partial_{\beta_j^\ast}W'=0$ is equivalent to the condition that $\sum_{ij} \beta_j\alpha_i\alpha_i^\ast=0$. Since the $\beta_j$ together form a monomorphism and the $\beta_j^\ast$ together form an epimorphism, this condition is equivalent to the condition $\sum_i \alpha_i\alpha_i^\ast\beta_j^\ast=0$ for each $j\in J$. But this is equivalent to the condition $\partial_{\beta_j}W=0$ in~\eqref{partial alpha-i W=0} using the substitutions~\eqref{proof_lemma2.11_1}  and~\eqref{proof_lemma2.11_2} above.

This shows that $Y$ is a representation of $J(Q',W')$ and $Y\in \,^{\perp}S_k'$. By naturality of the cokernel $Y_k$ in~\eqref{functorial exact sequence}, the assignment $Y=\psi_k(X)$ defines a functor $\psi_k:S_k^\perp\to \,^{\perp}S_k'$. In the opposite direction, for each $Y\in\,^{\perp}S_k'$, we let $X_k$ be the kernel of the epimorphism $(\beta_j^\ast):\bigoplus Y_j\to Y_k$. Let $\alpha_i:X_i=Y_i\to X_k$ be the unique morphism so that $\beta_{j'}\alpha_i=\gamma_{ij'}^\ast$ for $(i,j')\in P'$ and $\beta_j\alpha_i=-G_{ij}$ for $(i,j)\in P$. Finally, $\gamma_{ij}=\alpha_i^\ast\beta_j^\ast$ for all $(i,j)\in P$. The verification that $X\in S_k^\perp$ is analogous to the above discussion and it is clear that $Y\mapsto X$ gives the inverse of the functor $\psi_k$.
\end{proof}



\begin{lemm}\label{lem: rotation lemma for modules}
Using the same notation as in Lemma~\ref{lem: inductive step A=B for Q finite type}, let $M_1,\ldots,M_m\in S_k^\perp\subset \Lambda\text{-}\modr $. Then: 
\begin{enumerate}
\item\label{lemma2.12_1} %\emph{(1)} 
$\psi_k(M_1),\ldots,\psi_k(M_m)$ is a FHO sequence in $\Lambda'\text{-}\modr $ which lies in $\,^{\perp}S_k'$ if and only if 
\item\label{lemma2.12_2} %\emph{(2)} 
$S_k,M_1,\ldots,M_m$ is a FHO sequence in $\Lambda\text{-}\modr $.
\end{enumerate}
\end{lemm}

\begin{proof} Let $\cB=\cF(S_k\oplus M)$ where $M=M_1\oplus\cdots\oplus M_m$. Thus 
\[
	\cB=M^\perp\cap S_k^\perp= \{Y\in S_k^\perp\,|\, \Hom_{\Lambda}(M,Y)=0\}.
\]
We will show that both statements in the lemma are equivalent to the third statement:

\begin{enumerate}\setcounter{enumi}{2}
\item\label{proof_lemma2.12_3}  %(3) 
$M_1,\ldots,M_m$ is a maximal weakly FHO sequence in $S_k^\perp\cap\,^\perp \cB$.
\end{enumerate}


The equivalence $\eqref{lemma2.12_1}\Leftrightarrow\eqref{proof_lemma2.12_3}$ is clear. Since $\psi_k:S_k^\perp\cong \,^{\perp}S_k'$ by Lemma~\ref{lem: inductive step A=B for Q finite type}, we have
\[
	\cB':=\psi_k(\cB)=\psi_k(M)^\perp\cap \,^\perp S_k'=\{Y'\in\,^{\perp}S_k'\,|\, \Hom_{\Lambda'}(\psi_k(M),Y')=0\}.
\]
So, $\psi_k(S_k^\perp\cap\,^\perp \cB)=\,^{\perp}S_k'\cap\,^\perp\cB'=\cG(\psi_k(M))$ by Proposition~\ref{prop: HN stratification with 3 strata} with $S_k',\psi_k(M),\cB'$ playing the roles of $X,Y,\cC$. Since $\psi_k$ is an isomorphism, $M_1,\ldots,M_m$ is maximal weakly FHO in $S_k^\perp\cap\,^\perp\cB$ if and only if $\psi_k(M_1),\ldots,\psi_k(M_m)$ is maximal weakly FHO in $\psi_k(S_k^\perp\cap\,^\perp \cB)=\cG(\psi_k(M))$. By Definition~\ref{def of FHO} this is equivalent to~\eqref{lemma2.12_1}.


For the equivalence $\eqref{lemma2.12_2}\Leftrightarrow\eqref{proof_lemma2.12_3}$, note that~\eqref{proof_lemma2.12_3} is equivalent to the statement that $S_k$, $M_1,\ldots,M_m$ is weakly FHO in $\Lambda\text{-}\modr $ and that no objects of $\,^\perp\cB=\cG(S_k\oplus M)$ can be inserted between the $M_i$, after $M_m$ or before $M_1$ (and after $S_k$). Since these hold under condition~\eqref{lemma2.12_2} we have $\eqref{lemma2.12_2}\then \eqref{proof_lemma2.12_3}$. To show $\eqref{proof_lemma2.12_3}\then \eqref{lemma2.12_2}$ it remains to show one more condition: that no object of $\,^\perp\cB$ can be inserted before $S_k$ in the sequence.

Suppose not. Then there is a Schurian $\Lambda$-module $X\in\,^\perp\cB$ so that $\Hom_\Lambda(X,S_k\oplus M)=0$. Let $T\subsetneq X$ be the largest submodule having only $S_k$ in its composition series. Then $\Hom_\Lambda(S_k,X/T)=0$ and $\Hom_\Lambda(X/T,M)=0$. So, $S_k^\perp\cap\,^\perp\cB$ contains $X/T\neq0$. As in the proof of Proposition~\ref{prop: HN stratification with 3 strata}, any object $Z$ in $S_k^\perp\cap\,^\perp\cB$ of minimal length is Schurian. Then $Z,M_1,\ldots,M_m$ is a weakly FHO sequence in $S_k^\perp\cap\,^\perp \cB$ contradicting the maximality of the sequence $M_1,\ldots,M_m$. Therefore, all three statements are equivalent.
\end{proof}


